% !TeX program = pdflatex
\documentclass[11pt,a4paper]{amsart}

\newcommand{\notelanguage}{finnish}
% Shared preamble for course notes and exercise sets.

\usepackage[T1]{fontenc}
\usepackage{lmodern}
\providecommand{\notelanguage}{english}
\usepackage[finnish,english]{babel}
\AtBeginDocument{\expandafter\selectlanguage\expandafter{\notelanguage}}

\newcommand{\notesfinnishlanguage}{finnish}
\ifx\notelanguage\notesfinnishlanguage
\newcommand{\notestheoremname}{Lause}
\newcommand{\noteslemmaname}{Lemma}
\newcommand{\notespropositionname}{Propositio}
\newcommand{\notescorollaryname}{Korollaari}
\newcommand{\notesdefinitionname}{Määritelmä}
\newcommand{\notesexamplename}{Esimerkki}
\newcommand{\notesproblemname}{Tehtävä}
\newcommand{\notesexercisename}{Tehtävä}
\newcommand{\notesremarkname}{Huomautus}
\newcommand{\notessolutionname}{Ratkaisu}
\else
\newcommand{\notestheoremname}{Theorem}
\newcommand{\noteslemmaname}{Lemma}
\newcommand{\notespropositionname}{Proposition}
\newcommand{\notescorollaryname}{Corollary}
\newcommand{\notesdefinitionname}{Definition}
\newcommand{\notesexamplename}{Example}
\newcommand{\notesproblemname}{Problem}
\newcommand{\notesexercisename}{Exercise}
\newcommand{\notesremarkname}{Remark}
\newcommand{\notessolutionname}{Solution}
\fi

\usepackage[a4paper,margin=30mm]{geometry}
\usepackage{microtype}
\usepackage{graphicx}
\usepackage{enumitem}

\newcommand{\exerciseheading}[2]{%
  \par\noindent
  {\large\bfseries #1\par}
  \vspace{0.1em}
  \noindent{\large #2\par}
  \vspace{1.5em}
}

\usepackage{amsmath}
\usepackage{amssymb}
\usepackage{amsthm}
\usepackage{mathtools}
\usepackage{aliascnt}

% Row operations beside the affected row, e.g. \rowop{+2R_1}.
\newcommand{\rowop}[1]{\mathrel{\big|}\,#1}

% Linear systems: the argument gives the number of alignment pairs.
\newenvironment{linearsystem}[1]
{\left\{
  \begin{alignedat}{#1}}
    {
  \end{alignedat}\right.}

  \usepackage[hidelinks,pdfusetitle]{hyperref}

  \numberwithin{equation}{section}
  \setlist{itemsep=0.25em,topsep=0.5em}

  \theoremstyle{plain}
  \newtheorem{theorem}{\notestheoremname}[section]
  \newaliascnt{lemma}{theorem}
  \newtheorem{lemma}[lemma]{\noteslemmaname}
  \aliascntresetthe{lemma}
  \newaliascnt{proposition}{theorem}
  \newtheorem{proposition}[proposition]{\notespropositionname}
  \aliascntresetthe{proposition}
  \newaliascnt{corollary}{theorem}
  \newtheorem{corollary}[corollary]{\notescorollaryname}
  \aliascntresetthe{corollary}

  \theoremstyle{definition}
  \newaliascnt{definition}{theorem}
  \newtheorem{definition}[definition]{\notesdefinitionname}
  \aliascntresetthe{definition}
  \newaliascnt{example}{theorem}
  \newtheorem{example}[example]{\notesexamplename}
  \aliascntresetthe{example}
  \newaliascnt{problem}{theorem}
  \newtheorem{problem}[problem]{\notesproblemname}
  \aliascntresetthe{problem}
  \newtheorem{exercise}{\notesexercisename}

  \theoremstyle{remark}
  \newaliascnt{remark}{theorem}
  \newtheorem{remark}[remark]{\notesremarkname}
  \aliascntresetthe{remark}

  \newenvironment{solution}
  {
  \begin{proof}[\notessolutionname]}
    {
  \end{proof}}

  \usepackage[nameinlink,noabbrev]{cleveref}

  \crefname{theorem}{theorem}{theorems}
  \crefname{lemma}{lemma}{lemmas}
  \crefname{proposition}{proposition}{propositions}
  \crefname{corollary}{corollary}{corollaries}
  \crefname{definition}{definition}{definitions}
  \crefname{example}{example}{examples}
  \crefname{problem}{problem}{problems}
  \crefname{exercise}{exercise}{exercises}
  \crefname{remark}{remark}{remarks}

  \ifx\notelanguage\notesfinnishlanguage
  \crefname{theorem}{lause}{lauseet}
  \crefname{lemma}{lemma}{lemmat}
  \crefname{proposition}{propositio}{propositiot}
  \crefname{corollary}{korollaari}{korollaarit}
  \crefname{definition}{määritelmä}{määritelmät}
  \crefname{example}{esimerkki}{esimerkit}
  \crefname{problem}{tehtävä}{tehtävät}
  \crefname{exercise}{tehtävä}{tehtävät}
  \crefname{remark}{huomautus}{huomautukset}
  \fi

  \newcommand{\NN}{\mathbb{N}}
  \newcommand{\NNone}{\mathbb{N}_1}
  \newcommand{\ZZ}{\mathbb{Z}}
  \newcommand{\QQ}{\mathbb{Q}}
  \newcommand{\RR}{\mathbb{R}}
  \newcommand{\CC}{\mathbb{C}}
  \newcommand{\PP}{\mathbb{P}}

  \DeclareMathOperator{\im}{im}
  \DeclareMathOperator{\rank}{rank}
  \DeclareMathOperator{\syt}{syt}
  \DeclareMathOperator{\pyj}{pyj}
  \DeclareMathOperator{\dom}{dom}
  \DeclarePairedDelimiter{\abs}{\lvert}{\rvert}
  \DeclarePairedDelimiter{\norm}{\lVert}{\rVert}
  \DeclarePairedDelimiterXPP{\maxnorm}[1]{}{\lVert}{\rVert}{_\infty}{#1}
  \newcommand{\set}[1]{\left\lbrace#1\right\rbrace}
  \newcommand{\maxset}[1]{\max\set{#1}}
  \newcommand{\minset}[1]{\min\set{#1}}
  \newcommand{\powerset}[1]{\mathcal{P}\!\left(#1\right)}

\usepackage{pgfplots}
\pgfplotsset{compat=1.18}

\title{Lineaarialgebra ja matriisilaskenta I\\Harjoitus 5}
\author{}
\date{}

\begin{document}

\exerciseheading{Lineaarialgebra ja matriisilaskenta I}{Harjoitus 5}

\begin{exercise}
  Olkoon \(A\in\RR^{n\times n}\) kääntyvä matriisi. Osoita, että myös
  \(A^t\) on kääntyvä ja \((A^{-1})^t=(A^t)^{-1}\).
\end{exercise}

\begin{solution}
  Koska \(A\) on kääntyvä, sen käänteismatriisille pätee
  \(AA^{-1}=A^{-1}A=I\). Tulon transponointisäännön nojalla
  \[
    \begin{aligned}
      (A^{-1})^t A^t &= (AA^{-1})^t = I^t = I, \\
      A^t (A^{-1})^t &= (A^{-1}A)^t = I^t = I.
    \end{aligned}
  \]
  Siis \(A^t\) on kääntyvä ja sen käänteismatriisi on
  \((A^t)^{-1}=(A^{-1})^t\).
\end{solution}

\begin{exercise}
  Olkoon \(\theta\in\RR\). Osoita, että matriisi
  \[
    A=
    \begin{bmatrix}
      \cos\theta & -\sin\theta \\
      \sin\theta & \cos\theta
    \end{bmatrix}\in\RR^{2\times 2}
  \]
  on ortogonaalinen ja määritä \(A\):n käänteismatriisi \(A^{-1}\).
\end{exercise}

\begin{solution}
  Lasketaan
  \[
    \begin{aligned}
      A^tA
      &=
      \begin{bmatrix}
        \cos\theta & \sin\theta \\
        -\sin\theta & \cos\theta
      \end{bmatrix}
      \begin{bmatrix}
        \cos\theta & -\sin\theta \\
        \sin\theta & \cos\theta
      \end{bmatrix} \\
      &=
      \begin{bmatrix}
        \cos^2\theta+\sin^2\theta &
        -\cos\theta\sin\theta+\sin\theta\cos\theta \\
        -\sin\theta\cos\theta+\cos\theta\sin\theta &
        \sin^2\theta+\cos^2\theta
      \end{bmatrix}
      = I.
    \end{aligned}
  \]
  Siis \(A\) on ortogonaalinen. Korollaarin 4.5.4 nojalla
  \[
    A^{-1}=A^t=
    \begin{bmatrix}
      \cos\theta & \sin\theta \\
      -\sin\theta & \cos\theta
    \end{bmatrix}.\qedhere
  \]
\end{solution}

\begin{exercise}
  Etsi sellaiset matriisit \(A,B\in\RR^{2\times 2}\), että
  \(\det(A+B)\ne\det(A)+\det(B)\).
\end{exercise}

\begin{solution}
  Valitaan
  \[
    A=
    \begin{bmatrix}
      1 & 2 \\
      2 & 4
    \end{bmatrix},\qquad
    B=
    \begin{bmatrix}
      1 & 0 \\
      0 & 1
    \end{bmatrix}.
  \]
  Tällöin
  \[
    \det A=1\cdot 4-2\cdot 2=0,\qquad
    \det B=1\cdot 1-0\cdot 0=1.
  \]
  Siis
  \[
    \det(A+B)=
    \begin{vmatrix}
      2 & 2 \\
      2 & 5
    \end{vmatrix}
    =2\cdot 5-2\cdot 2=6\ne 0+1=\det A+\det B.\qedhere
  \]
\end{solution}

\begin{exercise}
  Olkoot
  \[
    A=
    \begin{bmatrix}
      a-3 & 5 \\
      -3 & a-2
    \end{bmatrix},\quad a\in\RR,\quad
    B=
    \begin{bmatrix}
      -2 & 7 & 6 \\
      5 & 1 & -2 \\
      3 & 8 & 4
    \end{bmatrix},\quad
    C=
    \begin{bmatrix}
      1 & -3 & 0 \\
      -2 & 4 & 1 \\
      5 & -2 & 2
    \end{bmatrix}.
  \]
  Laske determinantit \(\det A\), \(\det B\) ja \(\det C\).
  Ovatko matriisit \(A\), \(B\) ja \(C\) kääntyviä?
\end{exercise}

\begin{solution}
  Matriisin \(A\) determinantti on
  \[
    \det A=(a-3)(a-2)-5\cdot(-3)=a^2-5a+21.
  \]
  Yhtälön \(a^2-5a+21=0\) diskriminantti on
  \[
    (-5)^2-4\cdot 1\cdot 21=-59<0.
  \]
  Yhtälöllä ei siis ole reaalisia ratkaisuja. Siten \(\det A\ne 0\)
  kaikilla \(a\in\RR\).

  Kehitetään matriisin \(B\) determinantti ensimmäisen sarakkeen
  suhteen eli valitaan \(\ell=1\). Lauseen 5.9.4 nojalla
  \[
    \begin{aligned}
      \det B
      &= (-1)^2(-2)
      \begin{vmatrix}1 & -2 \\ 8 & 4
      \end{vmatrix}
      + (-1)^3\cdot 5
      \begin{vmatrix}7 & 6 \\ 8 & 4
      \end{vmatrix}
      + (-1)^4\cdot 3
      \begin{vmatrix}7 & 6 \\ 1 & -2
      \end{vmatrix} \\
      &= (-2)(4+16)-5(28-48)+3(-14-6) \\
      &= -40+100-60=0.
    \end{aligned}
  \]
  Kehitetään matriisin \(C\) determinantti kolmannen sarakkeen
  suhteen eli valitaan \(\ell=3\). Sarakkeen ensimmäinen alkio on
  nolla, joten Lauseen 5.9.4 nojalla
  \[
    \begin{aligned}
      \det C
      &= (-1)^5\cdot 1
      \begin{vmatrix}1 & -3 \\ 5 & -2
      \end{vmatrix}
      + (-1)^6\cdot 2
      \begin{vmatrix}1 & -3 \\ -2 & 4
      \end{vmatrix} \\
      &= -(-2+15)+2(4-6) \\
      &= -13-4=-17.
    \end{aligned}
  \]
  Lauseen 5.8.2 nojalla \(A\) on kääntyvä kaikilla \(a\in\RR\),
  \(C\) on kääntyvä ja \(B\) on singulaarinen.
\end{solution}

\begin{exercise}
  \begin{enumerate}[label=(\alph*), leftmargin=*]
    \item Olkoon \(A\in\RR^{n\times n}\). Osoita, että
      \(\det(kA)=k^n\det A\), missä \(k\in\RR\).
    \item Olkoot \(A,B\in\RR^{3\times 3}\) matriiseja, joille
      \(\det A=4\) ja \(\det B=5\). Määritä
      \[
        \text{(i)}\quad \det(AB),\qquad
        \text{(ii)}\quad \det(3A),\qquad
        \text{(iii)}\quad \det(2AB),\qquad
        \text{(iv)}\quad \det(A^{-1}B).
      \]
  \end{enumerate}
\end{exercise}

\begin{solution}
  \leavevmode\par
  \begin{enumerate}[label=(\alph*), leftmargin=*]
    \item Olkoon \(A_j\) matriisi, joka saadaan kertomalla \(A\):n
      \(j\) ensimmäistä saraketta luvulla \(k\). Osoitetaan
      induktiolla, että \(\det A_j=k^j\det A\) kaikilla
      \(1\le j\le n\).

      Kun \(j=1\), sarakelineaarisuuden (Lause 5.5.1) nojalla
      \(\det A_1=k\det A\).

      Oletetaan, että \(\det A_j=k^j\det A\) jollakin
      \(1\le j<n\). Matriisi \(A_{j+1}\) saadaan kertomalla
      \(A_j\):n \((j+1)\):s sarake luvulla \(k\). Sarakelineaarisuuden
      ja induktio-oletuksen nojalla
      \[
        \det A_{j+1}=k\det A_j=k(k^j\det A)=k^{j+1}\det A.
      \]
      Induktioperiaatteen nojalla väite pätee myös, kun \(j=n\).
      Koska \(A_n=kA\), saadaan \(\det(kA)=k^n\det A\).

    \item Kohdan (a) ja determinantin tulokaavan (Lause 5.6.1) nojalla
      \[
        \begin{aligned}
          \text{(i)}\quad \det(AB)
          &= \det A\det B=4\cdot 5=20, \\
          \text{(ii)}\quad \det(3A)
          &= 3^3\det A=27\cdot 4=108, \\
          \text{(iii)}\quad \det(2AB)
          &= 2^3\det(AB)=8\cdot 20=160, \\
          \text{(iv)}\quad \det(A^{-1}B)
          &= \frac{\det B}{\det A}=\frac54.
        \end{aligned}
      \]
      Viimeisessä kohdassa käytettiin yhtälöä
      \(\det A\det(A^{-1})=\det I=1\).
  \end{enumerate}
\end{solution}

\begin{exercise}
  Olkoon
  \[
    V=
    \begin{bmatrix}
      1 & x_1 & x_1^2 \\
      1 & x_2 & x_2^2 \\
      1 & x_3 & x_3^2
    \end{bmatrix}\in\RR^{3\times 3}
  \]
  \emph{Vandermonden matriisi}.
  \begin{enumerate}[label=(\alph*), leftmargin=*]
    \item Osoita, että \(\det V=(x_2-x_1)(x_3-x_1)(x_3-x_2)\).
    \item Mitkä ehdot lukujen \(x_1,x_2,x_3\) tulee täyttää, jotta
      \(V\) on kääntyvä?
  \end{enumerate}
\end{exercise}

\begin{solution}
  \leavevmode\par
  \begin{enumerate}[label=(\alph*), leftmargin=*]
    \item Kehitetään determinantti ensimmäisen sarakkeen suhteen
      eli valitaan \(\ell=1\). Ensimmäisen sarakkeen kaikki alkiot
      ovat ykkösiä, joten Lauseen 5.9.4 nojalla
      \[
        \begin{aligned}
          \det V
          &=
          \begin{vmatrix}x_2 & x_2^2 \\ x_3 & x_3^2
          \end{vmatrix}
          -
          \begin{vmatrix}x_1 & x_1^2 \\ x_3 & x_3^2
          \end{vmatrix}
          +
          \begin{vmatrix}x_1 & x_1^2 \\ x_2 & x_2^2
          \end{vmatrix} \\
          &= (x_2x_3^2-x_2^2x_3)-(x_1x_3^2-x_1^2x_3)
            +(x_1x_2^2-x_1^2x_2) \\
          &= x_2x_3^2-x_2^2x_3-x_1x_3^2+x_1^2x_3+x_1x_2^2-x_1^2x_2 \\
          &= (x_2-x_1)x_3^2-(x_2^2-x_1^2)x_3+x_1x_2(x_2-x_1) \\
          &= (x_2-x_1)\bigl[x_3^2-(x_1+x_2)x_3+x_1x_2\bigr] \\
          &= (x_2-x_1)\bigl[x_3(x_3-x_2)-x_1(x_3-x_2)\bigr] \\
          &= (x_2-x_1)(x_3-x_1)(x_3-x_2).
        \end{aligned}
      \]

    \item Lauseen 5.8.2 nojalla \(V\) on kääntyvä täsmälleen silloin,
      kun \(\det V\ne 0\). Kohdan (a) perusteella tämä pätee
      täsmälleen silloin, kun erotukset \(x_2-x_1\), \(x_3-x_1\)
      ja \(x_3-x_2\) ovat kaikki nollasta poikkeavia. Ehto voidaan
      kirjoittaa muodossa
      \[
        x_i\ne x_j
        \quad\text{kaikilla } i,j\in\{1,2,3\},\ \text{joilla } i\ne j.
        \qedhere
      \]
  \end{enumerate}
\end{solution}

\end{document}
